\documentclass[10pt,a4paper]{article}

%double si on veut une version prof et une version élève. simple si on veut une seule version.
\newcommand{\buildmode}{simple}

\InputIfFileExists{version.tex}{}{}
% Version (définie par le script)
\providecommand{\version}{eleve}

\usepackage{feuilleinterro}

%%%à modifier à chaque fois

\renewcommand{\niveau}{1M}
\renewcommand{\chapitre}{Proportions, pourcentages}
\renewcommand{\numchap}{0.2}
\renewcommand{\theme}{Statistiques}

\begin{document}

% =========================================================
% PAGE 1 : VERSION A
% =========================================================

\interroversion{A}

\textbf{Exercice 1 -- Calcul avec des fractions \hfill /3}

Calculer les expressions suivantes, en détaillant les étapes, et donner le résultat sous la forme d'une fraction irréductible.

\begin{enumerate}
\item \(A = \dfrac{6}{7}\times\dfrac{8}{9}\)

\vspace{2.5cm}

\item \(B = \dfrac{5}{6}-\dfrac{1}{3}\)

\vspace{2.5cm}
\end{enumerate}

\textbf{Exercice 2 -- Proportion, pourcentage et écriture décimale \hfill /3}

Dans un lycée de 200 élèves, 50 sont internes.

Quelle est la proportion d'élèves internes dans ce lycée ? Donner le résultat sous la forme d'un pourcentage, puis sous la forme d'un nombre décimal.

\vspace{3cm}

\textbf{Exercice 3 -- Tableau croisé \hfill /4}

Dans un lycée, on a répertorié les élèves selon leur sexe et leur régime (interne ou externe).

\begin{center}
\renewcommand{\arraystretch}{1.5}
\begin{tabular}{|l|c|c|c|}
\hline
 & \textbf{Interne} & \textbf{Externe} & \textbf{Total} \\
\hline
\textbf{Filles} & 50 & 30 & \\
\hline
\textbf{Garçons} & 70 & & 120 \\
\hline
\textbf{Total} & 120 & 80 & 200 \\
\hline
\end{tabular}
\end{center}

\begin{enumerate}
\item Compléter les deux cases manquantes du tableau. \hfill /2

\vspace{1.5cm}

\item Quelle est la proportion de filles parmi l'ensemble des élèves du lycée ? \hfill /1

\vspace{1.5cm}

\item Parmi les filles, quelle est la proportion d'élèves internes ? \hfill /1
\end{enumerate}

\newpage

% =========================================================
% PAGE 2 : VERSION B
% =========================================================

\interroversion{B}

\textbf{Exercice 1 -- Calcul avec des fractions \hfill /3}

Calculer les expressions suivantes, en détaillant les étapes, et donner le résultat sous la forme d'une fraction irréductible.

\begin{enumerate}
\item \(A = \dfrac{5}{9}\times\dfrac{7}{10}\)

\vspace{2.5cm}

\item \(B = \dfrac{7}{8}-\dfrac{1}{4}\)

\vspace{2.5cm}
\end{enumerate}

\textbf{Exercice 2 -- Proportion, pourcentage et écriture décimale \hfill /3}

Un club de sport compte 50 adhérents, dont 15 pratiquent la natation.

Quelle est la proportion d'adhérents pratiquant la natation ? Donner le résultat sous la forme d'un pourcentage, puis sous la forme d'un nombre décimal.

\vspace{3cm}

\textbf{Exercice 3 -- Tableau croisé \hfill /4}

Dans un club de sport, on a répertorié les adhérents selon leur sexe et le type d'abonnement souscrit (annuel ou mensuel).

\begin{center}
\renewcommand{\arraystretch}{1.5}
\begin{tabular}{|l|c|c|c|}
\hline
 & \textbf{Annuel} & \textbf{Mensuel} & \textbf{Total} \\
\hline
\textbf{Femmes} & 60 & 40 & \\
\hline
\textbf{Hommes} & 90 & & 150 \\
\hline
\textbf{Total} & 150 & 100 & 250 \\
\hline
\end{tabular}
\end{center}

\begin{enumerate}
\item Compléter les deux cases manquantes du tableau. \hfill /2

\vspace{1.5cm}

\item Quelle est la proportion de femmes parmi l'ensemble des adhérents ? \hfill /1

\vspace{1.5cm}

\item Parmi les femmes, quelle est la proportion d'abonnements annuels ? \hfill /1
\end{enumerate}

\newpage

% =========================================================
% PAGE 3 : VERSION C
% =========================================================

\interroversion{C}

\textbf{Exercice 1 -- Calcul avec des fractions \hfill /3}

Calculer les expressions suivantes, en détaillant les étapes, et donner le résultat sous la forme d'une fraction irréductible.

\begin{enumerate}
\item \(A = \dfrac{8}{11}\times\dfrac{5}{6}\)

\vspace{2.5cm}

\item \(B = \dfrac{7}{9}-\dfrac{1}{3}\)

\vspace{2.5cm}
\end{enumerate}

\textbf{Exercice 2 -- Proportion, pourcentage et écriture décimale \hfill /3}

Une entreprise emploie 200 salariés, dont 80 travaillent à temps partiel.

Quelle est la proportion de salariés à temps partiel dans cette entreprise ? Donner le résultat sous la forme d'un pourcentage, puis sous la forme d'un nombre décimal.

\vspace{3cm}

\textbf{Exercice 3 -- Tableau croisé \hfill /4}

Dans une entreprise, on a répertorié les salariés selon leur sexe et leur statut (cadre ou non-cadre).

\begin{center}
\renewcommand{\arraystretch}{1.5}
\begin{tabular}{|l|c|c|c|}
\hline
 & \textbf{Cadre} & \textbf{Non-cadre} & \textbf{Total} \\
\hline
\textbf{Femmes} & 50 & 50 & \\
\hline
\textbf{Hommes} & 70 & & 150 \\
\hline
\textbf{Total} & 120 & 130 & 250 \\
\hline
\end{tabular}
\end{center}

\begin{enumerate}
\item Compléter les deux cases manquantes du tableau. \hfill /2

\vspace{1.5cm}

\item Quelle est la proportion de femmes parmi l'ensemble des salariés ? \hfill /1

\vspace{1.5cm}

\item Parmi les femmes, quelle est la proportion de cadres ? \hfill /1
\end{enumerate}

\newpage

% =========================================================
% PAGE 4 : VERSION D
% =========================================================

\interroversion{D}

\textbf{Exercice 1 -- Calcul avec des fractions \hfill /3}

Calculer les expressions suivantes, en détaillant les étapes, et donner le résultat sous la forme d'une fraction irréductible.

\begin{enumerate}
\item \(A = \dfrac{9}{10}\times\dfrac{6}{7}\)

\vspace{2.5cm}

\item \(B = \dfrac{5}{8}-\dfrac{1}{4}\)

\vspace{2.5cm}
\end{enumerate}

\textbf{Exercice 2 -- Proportion, pourcentage et écriture décimale \hfill /3}

Un festival a accueilli 50 festivaliers, dont 35 avaient moins de 18 ans.

Quelle est la proportion de festivaliers de moins de 18 ans ? Donner le résultat sous la forme d'un pourcentage, puis sous la forme d'un nombre décimal.

\vspace{3cm}

\textbf{Exercice 3 -- Tableau croisé \hfill /4}

Dans un festival, on a répertorié les festivaliers selon le type de pass acheté et leur âge.

\begin{center}
\renewcommand{\arraystretch}{1.5}
\begin{tabular}{|l|c|c|c|}
\hline
 & \textbf{Moins de 18 ans} & \textbf{18 ans et plus} & \textbf{Total} \\
\hline
\textbf{Pass 3 jours} & 30 & 70 & \\
\hline
\textbf{Pass 1 jour} & 50 & & 150 \\
\hline
\textbf{Total} & 80 & 170 & 250 \\
\hline
\end{tabular}
\end{center}

\begin{enumerate}
\item Compléter les deux cases manquantes du tableau. \hfill /2

\vspace{1.5cm}

\item Quelle est la proportion de festivaliers ayant un pass 3 jours parmi l'ensemble des festivaliers ? \hfill /1

\vspace{1.5cm}

\item Parmi les festivaliers ayant un pass 3 jours, quelle est la proportion de moins de 18 ans ? \hfill /1
\end{enumerate}

\end{document}
